\section{Introduction}
\label{sec:intro}

% Background
% Status: reasoning model very strong and useful, many application
Recent advances in Large Language Model (LLM) reasoning have enabled broad applications across diverse domains~\citep{jaech2024openai-o1,guo2025deepseek-r1,yang2025qwen3}. 
With hundreds of billions of parameters, LLMs can generate complex Chain-of-Thought (CoT) to solve challenging tasks.
At the same time, smaller language models have also drawn increasing attention. 
With only a few billion parameters, they offer compelling alternatives: lower costs, faster inference, and suitability for edge computing~\citep{abdin2024phi3,minicpm4_2025,wang2025slm}.

% What problem
At this crossroads, enhancing reasoning capabilities under parameter constraints becomes a central challenge. 
A common approach is to distill smaller models to mimic LLM CoT trajectories.
% Why hard problem
However, not all tokens are equally predictable: certain tokens encode critical logic or reasoning directions that are fundamentally harder to predict~\citep{lin2024rho, fu2025r2r, wang2025beyond}.
With limited computation per output token, small models quickly hit a performance ceiling and mispredict some of these tokens.
Once critical errors occur, the reasoning trajectory can irrecoverably diverge, yielding drastically different outcomes.
% As pointed out by previous works, given the same context, while most next-token predictions between LLM and SLMs are the same same, a few hard and diverging tokens determines the subsequent reasoning trajectory, leading to completely different reasoning outcomes. 

% Previous work
Prior work proposes looped transformers to address this parameter–performance paradox~\citep{hutchins2022block, saunshi2025loopedTrans, zeng2025pondering, zhu2025ouro}. 
Instead of verbalizing the next token immediately after one forward pass, these models typically feed the last-layer hidden states back into the LLM for additional passes, refining representations in the latent space.
After certain iterations, the final hidden states pass through the language modeling head to generate the next token.
% \todo{use other model to be the motivation}
By uniformly scaling up iterations per token, these models can correct initially wrong token predictions, potentially increasing performance without increasing parameter count.

% Motivation: key insight
% However, we find that fixed-depth looped transformers sometimes degrade performance compared to direct response, primarily due to compute misallocation. As shown in Figure~\ref{fig:idea}, fine-tuning Qwen3-1.7B-Base to always perform two iterations per token yields no improvement over the direct-response baseline that iterates only once. We attribute this to \textit{latent overthinking}: similar to explicit CoT reasoning~\citep{wu2025morelesscot}, excessive iterations can revise correct answers into wrong ones. Since most cases, such as coherence or suffixes, are already predicted correctly at the first iteration, the opposite phenomenon, \textit{latent underthinking}, arises only on scarce tokens. We define tokens that cannot be accurately predicted in the current forward pass as \textit{hard} tokens.
% \zekai{In most cases, such as coherence or suffix prediction, the model predicts the correct token in the first iteration. However, a contrasting phenomenon, termed latent underthinking, occurs only with rare tokens. 
% We define tokens that cannot be accurately predicted in a single forward pass as hard tokens.}
% This motivates our central question:

% Motivation: key insight
However, we identify a \textit{latent overthinking} problem in looped transformers, where excessive iterations revise correct answers into wrong ones.
As shown in Figure~\ref{fig:landscape}, while the second iteration corrects 8.7\% of predictions, it also flips 2.1\% of correct ones into errors.
This occurs because most tokens, such as coherence or suffix tokens, are already predicted correctly after the first pass; further iterations may instead introduce harmful changes.
This mirrors overthinking in explicit CoT reasoning~\citep{wu2025morelesscot}, where additional reasoning steps degrade rather than improve answers.
% While the opposite \textit{latent underthinking} also exists for tokens that need more iterations to correctly predict, such cases are rarer.
% We define tokens that cannot be accurately predicted in a single forward pass as \textit{hard} tokens, and ask our central question:
This reveals a surprising opportunity:
\highlightblock{selectively skipping latent iterations on most tokens can further increase model performance.}
% If achieved, different iterations could specialize in distinct prediction focuses for better and more efficient reasoning.
% This divide-and-conquer scheme resembles Mixture-of-Experts (MoE), where iteration depths act as experts with complementary roles.
We validate this with an oracle policy that iterates only on initially mispredicted tokens, as shown in Table~\ref{tab:oracle_policy}. Compared to always iterating, this selective oracle can achieve up to 32\% higher accuracy with an optimized architecture.

% challenge
% Achieving dynamic latent iteration presents three main challenges. 
% First, the model architecture must support dynamic iteration depths, enabling cross-depth information flow while maintaining the parallel computation required for training and prefilling. It must also handle the shift in prediction objectives across iterations while maximizing parameter reuse. 
Achieving selective latent iteration presents three main challenges. 
First, the model architecture should enable cross-depth attention, allowing each iteration to access full context.
This is crucial because when early tokens skip deeper iterations, later tokens must still access their representations from shallower depths.
Meanwhile, this cross-depth flow cannot compromise the sequence-level parallelism essential for efficient training and prefilling.
Second, the model must adapt to distribution shifts across iterations, while maximizing parameter reuse.
Third, training must remain stable despite tightly coupled dependencies: the iteration policy depends on prediction quality at each depth, while that quality depends on the depths to which previous tokens are assigned by the iteration policy.

\begin{figure}[t]
    \centering
    \includegraphics[width=\columnwidth]{figure/idea.png}
    \caption{Latent iterations can fix wrong predictions, but can also \emph{overthink} and flip correct ones. \emph{Selective} iteration only when needed can improve reasoning with reduced computation.}   
    \label{fig:idea}
\end{figure}

To address these challenges, we propose \name, a looped transformer optimized for selective latent iteration.
As shown in Figure~\ref{fig:arch}, \name employs a neural decider to determine whether to iterate or verbalize each token.
We design a duo-causal attention mechanism to enable cross-depth attention with full sequence parallelism.
To specialize deeper iterations for current-token refinement and preserve strong first-pass predictions, we apply LoRA adapters solely at iterations $d>1$.
We train \name stably by aligning both the LLM backbone and iteration decider with a static oracle iteration policy.
We summarize our contributions as follows.

\begin{itemize}
    \item \xhdr{Selective Latent Iteration}
    We identify the latent overthinking phenomenon and quantify its influence on token prediction accuracy and downstream tasks.
    This insight motivates new directions where latent iteration is applied selectively to a few tokens for both better reasoning quality and efficiency.
    \item \xhdr{Specialized Model Architecture}
    We develop a model architecture that natively supports selective iterations.
    The dedicated duo-causal attention mechanism, LoRA adapters, and iteration decider enable efficient cross-depth information flow, objective transitions, and dynamic depth selection.
    \item \xhdr{Stable Training}
    We introduce a stable training scheme that uses a static oracle policy to decouple model adaptation and policy learning. It overcomes the circular dependency between iteration decisions and prediction quality.
\end{itemize}

We fine-tune \name from Qwen3-Base 0.6B and 1.7B on Open-R1, then test it across nine reasoning benchmarks spanning math, QA, and coding tasks. \name achieves average accuracy gains of 3.0\% and 3.8\% over standard single-iteration baselines.
Compared to looped transformer Ouro~\citep{zhu2025ouro} trained with same data, \name achieves 3.8-4.4\% gains while reducing latent iterations by 93\%.
